\setAuthor{Jonatan Kalmus}
\setRound{piirkonnavoor}
\setYear{2019}
\setNumber{G 9}
\setDifficulty{9}
\setTopic{Dünaamika}

\prob{Kiik}
Juku kiigub  nii suure hooga, et kiik jõuab haripunktis kiigepostide (kiige pöörlemistelje) kõrguseni. Suure hoo tõttu hakkab Juku muretsema, kas kiige postid ikka vastu peavad.\\
\osa Leida minimaalne ja maksimaalne jõud, mida Juku kiikumisel postidele avaldab.\\
\osa Leida minimaalne ja maksimaalne jõumoment, mida Juku kiikumisel postidele avaldab. Jõumoment arvutada kiigeposti alumise otsa ehk maasse kinnitamise punkti suhtes.
Kiige võll (pöörlemistelg) on kinnitatud kahe horisontaalse posti otsa kõrgusel $H$ maapinnast. Tasakaaluasendis on Juku massikeskme kõrgus maapinnast $h$. Juku mass on $m$ ning raskuskiirendus $g$. Kiige istme mass lugeda tühiseks. Lihtsustatult võib Jukut vaadelda punktmassina, mille kaugus kiige võllist ei muutu.\hint
Jukule mõjub kolm jõudu: vertikaalne raskuskiirendus, kiige nööridega paralleelne tsentrifugaaljõud ning nööride pinge, ehk teisisõnu Juku poolt postidedele avaldatav jõud. Kõik jõud on avaldatavad energia jäävuse seadusest. Jõumomente arvutades tuleb arvestada sellega, et kiigeposti mõjutab ainult kiige nööride pinge.\solu
\osa Minimaalne jõud: Haripunktis on Juku hetkeliselt vabalangemises ning mingit jõudu ei avaldu. Seega on minimaalne jõud: 
$$F_\mathrm{min} = 0.$$ 
Maksimaalne jõud: Kiikumise ajal rakenduvad postidele 2 jõudu, Juku raskusjõud ning pöörlemisest tulenev tsentrifugaaljõud. Mõlemad on maksimaalsed ning samasuunalised, kui kiik on vertikaalasendis. Olgu kiige pöörlemisraadius Juku massikeskme suhtes $R = H-h$. Kasutades energia jäävust kiige vertikaal- ja horisontaalasendi jaoks saame: 
$$mgR = \frac{mv^2}{2} \Rightarrow v^2 = 2gR.$$ 
Tsentrifugaaljõu valemist
$$F_C = m\frac{v^2}{R} = m\frac{2gR}{R} = 2mg.$$
Summarne jõud on seega
$$F_\mathrm{max} = F_C + F_R = 2mg + mg = 3mg.$$

\osa Minimaalne jõumoment: Kiige vertikaalasendis on jõuõlg kiigeposti alumise punkti suhtes $0$, seega minimaalne jõumoment on
$$\tau_\mathrm{min} = 0.$$ 
Maksimaalne jõumoment: Kuna võll kõrgusel $H$ on ainuke kiige kinnituspunkt, rakendub kogu Juku poolt rakendatav jõud sellesse punkti. Jõumomenti panustub ainult kiige kinnituspunkti rakenduva jõu horisontaalkomponent. Seega piisab maksimaalse horisontaalse jõu leidmisest. Tähistame kiige hetkeasukoha ja haripunkti tasandi vahelise kauguse $s$ ning nende vahelise nurga $\varphi$ (vt joonist). Tsentrifugaaljõud on alati suunatud piki raadiust ning sarnaselt eelnevaga saame tsentrifugaaljõu arvutamiseks rakendada energia jäävuse seadust:
$$F_C = m\frac{2gs}{R} = 2mg\sin{\varphi}.$$ 
Lisaks tuleb arvestada raskusjõu raadiusesihilise komponentiga $F_{R||} = F_R\sin{\varphi} = mg\sin{\varphi}$.
Seega rakendub piki raadiust võllile jõud
$$F_{||} = F_C + F_{R||} = 3mg\sin{\varphi}.$$ 
Meid huvitab selle jõu horisontaalkomponent
$$F_H = F_{||}\cos{\varphi} = 3mg\sin{\varphi}\cos{\varphi}.$$ 
Teades, et $\sin{2\varphi}=2\sin{\varphi}\cos{\varphi}$, saame
$$F_H = \frac{3}{2}mg\sin{2\varphi}.$$ 
Siinuse maksimaalne väärtus on $1$, mille saame $\varphi = \ang{45}$ korral, mis on antud ülesande jaoks saavutatav väärtus. Seega on maksimaalne horisontaalne jõud
$$F_H = \frac{3}{2}mg$$ 
ning sellele vastav maksimaalne jõumoment
$$\tau_\mathrm{max} = F_HH = \frac{3}{2}mgH.$$ 
\begin{center}
\includegraphics[width=0.35\linewidth]{2019-v2g-09-sol.png}
\end{center}\probend